
% =============================================================================
\chapter{Baseline Inicial Baseado em HTN}
% =============================================================================

\section{Motivação para iniciar pelo HTN}

A primeira etapa técnica do plano é construir um baseline HTN standalone. Essa
decisão foi tomada porque o HTN fornece a estrutura de alto nível necessária
para organizar o comportamento do agente.

Antes de integrar uma fonte de preferências e o seletor MORL de métodos, é necessário
garantir a representação de estados simbólicos, a decomposição de tarefas, a
verificação de pré-condições, a aplicação de efeitos, a geração e a execução de
planos, a atualização do estado por sensores e o replanejamento quando o plano
deixa de ser válido.

\section{Componentes implementados}

A implementação atual contém uma classe de estado de mundo e representações de
pré-condições, efeitos, tarefas primitivas, tarefas compostas e métodos de
decomposição. Sobre essas abstrações, o domínio HTN e o \textit{planner} recursivo
produzem planos, enquanto uma interface abstrata de ação e seus estados de
execução permitem ao agente executor realizar as tarefas. O sistema também
inclui sensores, um mecanismo de notificação de mudança de estado, uma abstração
de busca de caminhos com BFS para o grid, o ambiente GridWorld e um
visualizador em terminal.

\section{Estado simbólico}

O estado simbólico representa fatos observáveis sobre o mundo. Ele funciona
como uma estrutura de pares chave-valor. No GridWorld atual, inclui
\texttt{agent\_x}, \texttt{agent\_y}, \texttt{has\_key},
\texttt{door\_open}, \texttt{at\_key}, \texttt{at\_door},
\texttt{at\_goal} e \texttt{done}.

O estado pode ser copiado.
Isso é essencial porque o planner precisa simular efeitos sem modificar o
ambiente real.

\section{Tarefas primitivas}

Uma tarefa primitiva representa uma unidade planejável que pode ser executada
por uma ação concreta.
Formalmente, uma tarefa primitiva pode ser descrita por:

\begin{equation}
t_p = \langle nome, acao, pre, eff \rangle.
\end{equation}

Aqui, $pre$ representa o conjunto de pré-condições e $eff$ representa o
conjunto de efeitos.

\section{Tarefas compostas}

Uma tarefa composta representa uma intenção de alto nível.
Ela não é executada diretamente.
Em vez disso, é decomposta em subtarefas por meio de métodos.

Formalmente:

\begin{equation}
t_c = \langle nome, metodos \rangle.
\end{equation}

\section{Métodos}

Um método representa uma forma possível de decompor uma tarefa composta.
Ele possui pré-condições e uma lista de subtarefas:

\begin{equation}
m = \langle pre(m), subtasks(m) \rangle.
\end{equation}

Um método só pode ser usado quando suas pré-condições são satisfeitas.

\section{Plano}

Um plano é uma sequência de tarefas primitivas:

\begin{equation}
\pi = \langle a_1,a_2,\dots,a_n \rangle.
\end{equation}

Um plano é válido quando cada ação é aplicável no estado produzido pela
simulação das ações anteriores.

\section{Planner recursivo}

O método \texttt{build\_plan} percorre as tarefas raiz do domínio usando uma
cópia do estado observado. Quando a expansão de uma tarefa tem sucesso, suas
tarefas primitivas são anexadas ao plano e o estado simulado resultante é
propagado para a próxima tarefa raiz, como mostra o Algoritmo
\ref{alg:construir-plano}.

\begin{algorithm}[H]
\caption{Método \texttt{build\_plan} do baseline HTN}
\label{alg:construir-plano}
\begin{algorithmic}[1]
\Require domínio $D$, estado observado $WS_0$
\Ensure plano de tarefas primitivas $P$
\State $P \gets \langle\rangle$; $WS \gets \mathrm{Copia}(WS_0)$
\ForAll{$\tau \in D.\mathrm{tasks}$}
  \State $r \gets \Call{PlanejarRecursivamente}{\langle\rangle, WS, \tau}$
  \If{$r \neq \mathrm{Falha}$}
    \State $(P_\tau, WS) \gets r$
    \State $P \gets P \mathbin{\Vert} P_\tau$
  \EndIf
\EndFor
\State \Return $P$
\end{algorithmic}
\end{algorithm}

\section{Planejamento recursivo e backtracking}

O método \texttt{recursive\_planning} expande uma tarefa sobre um estado
simulado. Para uma tarefa primitiva, ele verifica pré-condições, inclui a tarefa
no plano do ramo e aplica seus efeitos apenas à cópia de estado. Para uma tarefa
composta, obtém os métodos aplicáveis e os tenta na ordem definida no domínio.
Se alguma subtarefa falhar, o método inteiro é descartado e o próximo método é
tentado. Esse é o backtracking em nível de método implementado no baseline.

\begin{algorithm}[H]
\caption{Método \texttt{recursive\_planning} com backtracking por método}
\label{alg:planejar-recursivamente}
\begin{algorithmic}[1]
\Require prefixo $P$, estado simulado $WS$, tarefa $\tau$
\Ensure $(P', WS')$ ou $\mathrm{Falha}$
\If{$\tau$ é primitiva}
  \If{$\neg\mathrm{PreCondicoesSatisfeitas}(\tau, WS)$}
    \State \Return $\mathrm{Falha}$
  \EndIf
  \State $P' \gets \mathrm{Copia}(P)$; $WS' \gets \mathrm{Copia}(WS)$
  \State $P' \gets P' \mathbin{\Vert} \langle\tau\rangle$
  \State $WS' \gets \mathrm{AplicarEfeitos}(\tau, WS')$
  \State \Return $(P', WS')$
\EndIf
\If{$\tau$ é composta}
  \ForAll{$m \in \mathrm{MetodosAplicaveis}(\tau, WS)$, na ordem do domínio}
    \State $P_m \gets \mathrm{Copia}(P)$; $WS_m \gets \mathrm{Copia}(WS)$
    \State $falhou \gets \mathrm{falso}$
    \ForAll{$\tau' \in m.\mathrm{subtasks}$}
      \State $r \gets \Call{PlanejarRecursivamente}{P_m, WS_m, \tau'}$
      \If{$r = \mathrm{Falha}$}
        \State $falhou \gets \mathrm{verdadeiro}$; \State \textbf{break}
      \EndIf
      \State $(P_m, WS_m) \gets r$
    \EndFor
    \If{$\neg falhou$}
      \State \Return $(P_m, WS_m)$
    \EndIf
  \EndFor
\EndIf
\State \Return $\mathrm{Falha}$
\end{algorithmic}
\end{algorithm}

% =============================================================================
\chapter{Ambiente Experimental GridWorld}
% =============================================================================

\section{Motivação}

O GridWorld foi escolhido como ambiente inicial por ser simples, discreto,
controlável e adequado para validar planejamento simbólico.
Nesta etapa, o objetivo não é treinar uma política de RL, mas validar a
arquitetura do baseline HTN.

\section{Modelo formal do ambiente}

O ambiente é modelado como um sistema de transição determinístico:

\begin{equation}
\mathcal{E} = \langle S,A,T,R,s_0,G \rangle.
\end{equation}

Nessa formulação, $S$ é o conjunto de estados possíveis, $A$ é o conjunto de
ações discretas, $T$ é a função de transição, $R$ é a recompensa, $s_0$ é o
estado inicial e $G$ é a condição de objetivo. Na versão atual, $R=1$ somente
quando o objetivo é atingido e $R=0$ nos demais passos. O ambiente não fornece
recompensa vetorial nesta etapa; seu papel é validar o ciclo HTN em um cenário
controlado.

\section{Estado do ambiente}

O estado concreto do GridWorld inclui:

\begin{equation}
s =
\langle
p_a,
p_k,
p_d,
p_g,
O,
hasKey,
doorOpen,
done
\rangle.
\end{equation}

Nessa representação, $p_a$, $p_k$, $p_d$ e $p_g$ são, respectivamente, as
posições do agente, da chave, da porta e do objetivo; $O$ é o conjunto de
obstáculos; e $hasKey$, $doorOpen$ e $done$ registram a posse da chave, a
abertura da porta e o término do episódio. A observação retornada pelo ambiente
preserva esses fatos concretos; o sensor os traduz para os predicados usados
pelo domínio HTN.

\section{Ações do ambiente}

O espaço de ações é discreto, com seis operações: mover para cima, direita,
baixo ou esquerda, pegar a chave e abrir a porta. Movimentos que saem do grid,
atingem um obstáculo ou tentam entrar no objetivo com a porta fechada não
alteram a posição. A coleta da chave só é efetiva sobre sua célula, e a abertura
da porta exige que o agente esteja em sua célula e já possua a chave.

O Algoritmo \ref{alg:gridworld-step} descreve o método \texttt{step}, que
processa uma ação, atualiza o término do episódio e retorna o contrato de cinco
valores do Gymnasium.

\begin{algorithm}[H]
\caption{Método \texttt{step} do ambiente GridWorld}
\label{alg:gridworld-step}
\begin{algorithmic}[1]
\Require ação discreta $a$
\Ensure observação $o$, recompensa $r$, término $terminated$, truncamento $truncated$, informação $info$
\If{$done$}
  \State \Return $(\mathrm{Obs}(), 0, \mathrm{verdadeiro}, \mathrm{falso}, \{\})$
\EndIf
\If{$a \in \{\mathrm{cima},\mathrm{direita},\mathrm{baixo},\mathrm{esquerda}\}$}
  \State \Call{TentarMover}{$a$}
\ElsIf{$a = \mathrm{pegarChave}$}
  \State \Call{TentarPegarChave}{}
\ElsIf{$a = \mathrm{abrirPorta}$}
  \State \Call{TentarAbrirPorta}{}
\Else
  \State sinalizar ação inválida
\EndIf
\State $done \gets (p_a=p_g) \land doorOpen$
\State $r \gets 1$ se $done$; caso contrário, $r \gets 0$
\State \Return $(\mathrm{Obs}(), r, done, \mathrm{falso}, \{\})$
\end{algorithmic}
\end{algorithm}

\section{Objetivo do ambiente}

O objetivo é alcançar a posição final com a porta aberta.
Portanto, o episódio termina quando:

\begin{equation}
p_a = p_g
\quad \text{e} \quad
doorOpen = true.
\end{equation}

\section{Configuração}

\texttt{GridWorldConfig} define largura, altura, posições do agente, chave,
porta e objetivo, obstáculos fixos e aleatórios, além dos estados iniciais de
posse da chave e abertura da porta. Uma posição definida como \texttt{None} é
sorteada no reinício do episódio sem sobrepor entidades especiais ou obstáculos.
Assim, posições fixas permitem testes determinísticos, enquanto uma
\textit{seed} torna cenários aleatórios reproduzíveis.

O Algoritmo \ref{alg:gridworld-reset} descreve o método \texttt{reset}.
Ele resolve o layout de cada episódio, restaura os estados iniciais configurados
e produz a primeira observação.

\begin{algorithm}[H]
\caption{Método \texttt{reset} do ambiente GridWorld}
\label{alg:gridworld-reset}
\begin{algorithmic}[1]
\Require semente opcional $seed$
\Ensure observação inicial $o$ e informação $info$
\State inicializar o gerador aleatório com $seed$
\State resolver posições fixas e aleatórias; adicionar obstáculos aleatórios
\State $p_a \gets p_{\mathrm{inicio}}$
\State $hasKey \gets initialHasKey$; $doorOpen \gets initialDoorOpen$
\State $done \gets (p_a=p_g) \land doorOpen$
\State \Return $(\mathrm{Obs}(), \{\})$
\end{algorithmic}
\end{algorithm}

% =============================================================================
\chapter{Domínio HTN do GridWorld}
% =============================================================================

\section{Especificação do domínio HTN}

O capítulo anterior descreve a dinâmica concreta do GridWorld; este capítulo
especifica como essa dinâmica é abstraída em tarefas e métodos HTN. A tarefa
raiz é \texttt{escape\_grid}. As folhas primitivas são
\texttt{go\_to\_key}, \texttt{pickup\_key}, \texttt{go\_to\_door},
\texttt{open\_door} e \texttt{go\_to\_goal}; elas associam uma ação concreta
a pré-condições e efeitos simbólicos. Em particular, as tarefas de navegação
atualizam \texttt{agent\_x}, \texttt{agent\_y} e os predicados
\texttt{at\_key}, \texttt{at\_door} ou \texttt{at\_goal};
\texttt{pickup\_key} produz \texttt{has\_key=true};
\texttt{open\_door} produz \texttt{door\_open=true}; e
\texttt{go\_to\_goal} produz \texttt{done=true}.

As tarefas compostas \texttt{ensure\_has\_key} e
\texttt{ensure\_door\_open} expressam, respectivamente, a obtenção da chave e
a abertura da porta. Cada uma tem um método vazio quando a condição já está
satisfeita, evitando ações artificiais. A tarefa \texttt{escape\_grid} leva
diretamente ao objetivo quando a porta está aberta; caso contrário, primeiro
garante a chave e a abertura da porta. A Figura
\ref{fig:dominio-htn-gridworld} sintetiza essa construção a partir do layout efetivo
do ambiente.

\begin{figure}[H]
\centering
\resizebox{0.92\textwidth}{!}{%
\begin{tikzpicture}[
    font=\footnotesize,
    >=Latex,
    task/.style={rectangle, rounded corners=2mm, draw=htnblue, thick,
        fill=htnblue!8, align=center, minimum width=32mm, minimum height=7mm},
    primitive/.style={rectangle, rounded corners=2mm, draw=gamegreen, thick,
        fill=gamegreen!8, align=center, minimum width=24mm, minimum height=7mm},
    noop/.style={rectangle, rounded corners=2mm, draw=envgray, thick,
        fill=envgray!7, align=center, minimum width=21mm, minimum height=7mm},
    condition/.style={font=\scriptsize, fill=white, inner sep=1.5pt, align=center},
    arrow/.style={->, thick}
]
    % Decomposição da tarefa raiz.
    \node[task] (escape) at (0,0) {\texttt{escape\_grid}};
    \node[primitive] (direct-goal) at (-5.0,-2.1) {\texttt{go\_to\_goal}};
    \node[task] (key) at (2.4,-2.1) {\texttt{ensure\_has\_key}};
    \node[task] (door) at (2.4,-4.8) {\texttt{ensure\_door\_open}};
    \node[primitive] (final-goal) at (2.4,-7.5) {\texttt{go\_to\_goal}};

    \draw[arrow] (escape) -- node[condition, above left, pos=0.57] {\texttt{door\_open}} (direct-goal);
    \draw[arrow] (escape) -- node[condition, above right, pos=0.57] {$\neg$\texttt{door\_open}} (key);
    \draw[arrow] (key) -- (door);
    \draw[arrow] (door) -- (final-goal);

    % Métodos de ensure_has_key, em uma faixa exclusiva.
    \node[noop] (key-noop) at (-0.2,-3.7) {método vazio};
    \node[primitive] (go-key) at (5.0,-3.7) {\texttt{go\_to\_key}};
    \node[primitive] (pick-key) at (8.3,-3.7) {\texttt{pickup\_key}};
    \draw[arrow] (key) -- node[condition, pos=.55, above left] {\texttt{has\_key}} (key-noop);
    \draw[arrow] (key) -- node[condition, pos=.55, above right] {$\neg$\texttt{has\_key}} (go-key);
    \draw[arrow] (go-key) -- (pick-key);

    % Métodos de ensure_door_open, em uma segunda faixa exclusiva.
    \node[noop] (door-noop) at (-0.2,-6.4) {método vazio};
    \node[primitive] (go-door) at (5.0,-6.4) {\texttt{go\_to\_door}};
    \node[primitive] (open-door) at (8.3,-6.4) {\texttt{open\_door}};
    \draw[arrow] (door) -- node[condition, pos=0.66, above left] {\texttt{door\_open}} (door-noop);
    \draw[arrow] (door) -- node[condition, above right, pos=0.67] {$\neg$\texttt{door\_open} $\land$ \texttt{has\_key}} (go-door);
    \draw[arrow] (go-door) -- (open-door);
\end{tikzpicture}
}
\caption{Decomposições do domínio HTN do GridWorld. Nós azuis são tarefas
compostas, nós verdes são tarefas primitivas e nós cinza representam métodos
vazios para subobjetivos já satisfeitos.}
\label{fig:dominio-htn-gridworld}
\end{figure}


Por exemplo, quando o agente não possui a chave e a porta está fechada, a
decomposição produz o plano
\begin{equation}
\begin{aligned}
\pi = \langle {} &\texttt{go\_to\_key},\texttt{pickup\_key},
\texttt{go\_to\_door},\\
&\texttt{open\_door},\texttt{go\_to\_goal} \rangle.
\end{aligned}
\end{equation}
Se a porta já estiver aberta, o método correspondente reduz o plano à navegação
até o objetivo. Essa redução é consequência das pré-condições e dos métodos
vazios do domínio, e não de uma regra especial no ambiente.

% =============================================================================
\chapter{Execução por Agente}
% =============================================================================

\section{Separação entre planejamento e execução}

Uma decisão arquitetural importante foi separar o planner do agente executor.
O planner decompõe tarefas, verifica pré-condições, aplica efeitos apenas ao
estado simulado e retorna folhas primitivas. O agente, por sua vez, mantém o
plano em execução, solicita um novo plano quando necessário, executa a tarefa
corrente e valida o plano remanescente após mudanças observadas no mundo.

\section{Execução por ticks}

A execução ocorre em ciclos discretos chamados ticks.
Em cada tick, o agente pode reconstruir o plano, executar uma etapa da tarefa
corrente e atualizar o plano pendente. Isso é necessário porque ações de
navegação podem durar vários ticks. O Algoritmo \ref{alg:agent-tick} descreve
o método \texttt{tick} atualmente implementado.

\begin{algorithm}[H]
\caption{Método \texttt{tick} do agente HTN}
\label{alg:agent-tick}
\begin{algorithmic}[1]
\Require mundo de execução $W$
\Ensure resultado do tick $R$
\State $replanned \gets \mathrm{falso}$; $planned \gets \langle\rangle$
\If{\Call{DeveReplanejar}{}}
  \State $plan \gets \Call{build\_plan}{}$; $worldStateChanged \gets \mathrm{falso}$
  \State $replanned \gets \mathrm{verdadeiro}$; $planned \gets \Call{NomesDoPlano}{}$
  \If{$plan=\langle\rangle$}
    \State \Return resultado sem tarefa e mensagem ``nenhum plano válido''
  \EndIf
\EndIf
\State $\tau \gets plan[0]$
\If{$\tau$ não é primitiva}
  \State remover $\tau$ do plano; \Return resultado de tarefa ignorada
\EndIf
\State $status \gets \Call{Executar}{\tau.\mathrm{action},W}$
\If{$status=\mathrm{SUCCESS}$}
  \State remover $\tau$ do plano
\ElsIf{$status=\mathrm{FAILURE}$}
  \State $plan \gets \langle\rangle$
\ElsIf{$status \neq \mathrm{RUNNING}$}
  \State sinalizar estado de ação desconhecido
\EndIf
\State \Return $(\tau.\mathrm{name},status,replanned,planned,\Call{NomesDoPlano}{})$
\end{algorithmic}
\end{algorithm}

\section{Status das ações}

Cada ação retorna \texttt{RUNNING} enquanto permanece em execução,
\texttt{SUCCESS} quando conclui a tarefa ou \texttt{FAILURE} quando não pode
concluí-la. Na navegação, por exemplo, o status permanece \texttt{RUNNING} até
o agente alcançar o destino; a ausência de rota resulta em \texttt{FAILURE}.

\section{Resultado de um tick}

O objeto \texttt{AgentTickResult} registra a tarefa executada, o status
retornado, a ocorrência de replanejamento, os nomes das tarefas do plano
gerado e do plano restante, além de uma mensagem opcional. Essas informações
permitem depurar e analisar o ciclo sem acoplar o agente ao visualizador.

% =============================================================================
\chapter{Sensores e Replanejamento}
% =============================================================================

\section{Motivação dos sensores}

O planner não deve acessar diretamente o ambiente concreto.
Em vez disso, ele deve operar sobre um estado simbólico atualizado por sensores.

Um sensor pode ser representado como uma função:

\begin{equation}
\sigma_i:E \rightarrow O_i.
\end{equation}

Aqui, $E$ representa o ambiente e $O_i$ representa uma observação produzida
pelo sensor.

\section{Sensor do GridWorld}

O \texttt{GridWorldSensor} atualiza no \texttt{WorldState} as coordenadas do
agente, da chave, da porta e do objetivo, os fatos \texttt{has\_key},
\texttt{door\_open} e \texttt{done}, e os predicados derivados
\texttt{at\_key}, \texttt{at\_door} e \texttt{at\_goal}. Dessa forma, o
domínio HTN recebe fatos simbólicos sem depender da representação concreta do
ambiente Gymnasium.

\section{Delegate de mudança de estado}

Depois que os sensores atualizam o estado, o \texttt{SensorSystem} notifica os
assinantes por meio de um \textit{delegate}. O agente copia o estado observado,
atualiza a cópia do planner e marca que houve mudança de mundo, como mostrado
na Figura \ref{fig:sensores-replanejamento}.

\begin{figure}[H]
\centering
\begin{tikzpicture}[
    font=\small,
    >=Latex,
    node distance=7mm and 15mm,
    block/.style={rectangle, rounded corners=2mm, draw, thick, align=center,
        text width=42mm, minimum height=9mm},
    runtime/.style={block, draw=gamegreen, fill=gamegreen!8},
    symbolic/.style={block, draw=htnblue, fill=htnblue!8},
    control/.style={block, draw=fallbackorange, fill=fallbackorange!10},
    decision/.style={diamond, aspect=2.1, draw=black!70, thick, align=center,
        fill=black!3, text width=29mm, inner sep=1.5pt},
    arrow/.style={->, thick}
]
    \node[runtime] (env) {Ambiente GridWorld};
    \node[runtime, below=of env] (sensor) {Sensor atualiza fatos simbólicos};
    \node[symbolic, below=of sensor] (state) {\texttt{WorldState} atualizado};
    \node[control, below=of state] (update) {\textit{Delegate} notifica o agente, que atualiza o planner};
    \node[control, below=of update] (validate) {Validar o plano remanescente no próximo tick};
    \node[decision, below=of validate] (valid) {O plano continua válido?};
    \node[symbolic, below left=of valid] (keep) {Manter execução do plano};
    \node[control, below right=of valid] (replan) {Solicitar novo plano HTN};

    \draw[arrow, gamegreen] (env) -- (sensor);
    \draw[arrow, htnblue] (sensor) -- (state);
    \draw[arrow, fallbackorange] (state) -- (update);
    \draw[arrow, fallbackorange] (update) -- (validate);
    \draw[arrow] (validate) -- (valid);
    \draw[arrow, htnblue] (valid) -- node[above left] {sim} (keep);
    \draw[arrow, fallbackorange] (valid) -- node[above right] {não} (replan);
\end{tikzpicture}
\caption{Atualização sensorial e replanejamento preguiçoso. Uma mudança de
estado não descarta imediatamente o plano: o agente valida o plano remanescente
no tick seguinte e só solicita novo planejamento se ele se tornou inválido.}
\label{fig:sensores-replanejamento}
\end{figure}


\section{Validação preguiçosa}

O agente não descarta o plano imediatamente após qualquer mudança de estado.
Em vez disso, valida se o plano restante ainda pode ser executado.
Essa estratégia evita replanejamento desnecessário.

Dado um plano restante:

\begin{equation}
\pi_r =
\langle a_i,a_{i+1},\dots,a_n \rangle,
\end{equation}

o agente simula as tarefas restantes a partir do estado atual. Se todas as
pré-condições forem satisfeitas, o plano continua válido. O Algoritmo
\ref{alg:validar-plano} explicita essa validação, incluindo os efeitos de cada
tarefa primitiva sobre a cópia de estado.

\begin{algorithm}[H]
\caption{Método \texttt{\_is\_plan\_still\_valid}}
\label{alg:validar-plano}
\begin{algorithmic}[1]
\Require plano restante $\pi_r$, estado observado $WS$
\Ensure \texttt{verdadeiro} se $\pi_r$ permanece executável
\State $WS' \gets \mathrm{Copia}(WS)$
\ForAll{$\tau \in \pi_r$}
  \If{$\tau$ não é primitiva}
    \State \textbf{continue}
  \EndIf
  \If{$\neg\mathrm{PreCondicoesSatisfeitas}(\tau,WS')$}
    \State \Return \texttt{falso}
  \EndIf
  \State $WS' \gets \mathrm{AplicarEfeitos}(\tau,WS')$
\EndFor
\State \Return \texttt{verdadeiro}
\end{algorithmic}
\end{algorithm}

\section{Condição de replanejamento}

O replanejamento ocorre quando não existe plano ou quando uma mudança observada
invalida o plano remanescente. Uma falha de ação limpa o plano e, por isso,
provoca replanejamento no tick seguinte. Formalmente:

\begin{equation}
\mathrm{Replanejar}
\iff
(\pi_r=\langle\rangle)
\lor
(\mathrm{mudancaDeMundo}\land\neg\mathrm{Valido}(\pi_r,WS)).
\end{equation}


Esse mecanismo prepara a arquitetura para ambientes mais dinâmicos, mesmo que
o GridWorld inicial seja determinístico.

% =============================================================================
\chapter{Pathfinding e Navegação}
% =============================================================================

\section{Separação entre intenção e movimento}

O HTN define intenções simbólicas, como ir até a chave, ir até a porta ou ir
até o objetivo.
O pathfinder transforma essas intenções em movimentos concretos.

Essa separação evita acoplamento entre o domínio HTN e o algoritmo de
navegação.

\section{Grafo do GridWorld}

O GridWorld pode ser modelado como um grafo:

\begin{equation}
G_{grid} = \langle V,E \rangle.
\end{equation}

O conjunto $V$ contém células livres do grid.
O conjunto $E$ contém conexões entre células ortogonalmente adjacentes.

Uma célula é válida se está dentro dos limites do grid e não pertence ao
conjunto de obstáculos.

\section{Busca em largura}

A implementação atual usa busca em largura, ou BFS. Como as arestas do grid
possuem custo uniforme, BFS encontra o menor caminho em número de passos.

O Algoritmo \ref{alg:bfs-grid} corresponde ao método
\texttt{GridPathfinder.find\_path}. Ele expande vizinhos ortogonais na ordem
cima, direita, baixo e esquerda, ignora células bloqueadas ou já visitadas e
reconstrói o caminho a partir do mapa de predecessores.

\begin{algorithm}[H]
\caption{Método \texttt{find\_path} por busca em largura}
\label{alg:bfs-grid}
\begin{algorithmic}[1]
\Require início $s$, objetivo $g$, contexto $(largura,altura,bloqueadas)$
\Ensure caminho de $s$ até $g$ ou lista vazia
\If{$s=g$}
  \State \Return $[s]$
\EndIf
\State $fronteira \gets \mathrm{Fila}([s])$; $predecessor[s] \gets \emptyset$
\While{$fronteira$ não está vazia}
  \State $atual \gets \Call{RemoverInicio}{fronteira}$
  \If{$atual=g$}
    \State \textbf{break}
  \EndIf
  \ForAll{$vizinho$ ortogonal válido de $atual$}
    \If{$vizinho \in bloqueadas$ \textbf{ ou } $vizinho \in predecessor$}
      \State \textbf{continue}
    \EndIf
    \State $predecessor[vizinho] \gets atual$
    \State \Call{InserirFim}{fronteira,vizinho}
  \EndFor
\EndWhile
\If{$g \notin predecessor$}
  \State \Return $[\ ]$
\EndIf
\State \Return \Call{ReconstruirCaminho}{$predecessor,g$}
\end{algorithmic}
\end{algorithm}

A ação \texttt{NavigateToPositionAction} consulta esse caminho a cada tick,
executa apenas o próximo movimento por meio de \texttt{env.step} e retorna
\texttt{SUCCESS} ao alcançar o alvo. Se o caminho tiver menos de duas células
antes de alcançar o alvo, a ação retorna \texttt{FAILURE}; caso contrário,
retorna \texttt{RUNNING}. A recomputação do caminho torna a navegação reativa
ao estado corrente do ambiente.

% =============================================================================
\chapter{Evolução para o Experimento Multiobjetivo}
% =============================================================================

\section{Papel do baseline atual}

O GridWorld existente permanece útil como prova de conceito do mecanismo HTN,
da execução por ticks, dos sensores, da navegação e do replanejamento.
Ele não constitui, isoladamente, a avaliação final da hipótese.
Sua função é fornecer uma base testável sobre a qual a interface de seleção de
métodos possa ser introduzida sem misturar falhas do planner com falhas de
aprendizado.
Ele ainda não oferece alternativas genuinamente conflitantes, recompensa
vetorial, nem um protocolo de preferência-condicionada; portanto, não é por si
só um benchmark MORL adequado.

\section{Alteração necessária no planner}

O planner atual percorre métodos aplicáveis segundo uma ordem definida no
domínio e usa backtracking quando uma decomposição falha.
Para a arquitetura proposta, a etapa de seleção deverá ser desacoplada da
enumeração fixa.
No caminho online normal, o seletor MORL recebe a máscara dos métodos aplicáveis,
faz uma inferência e escolhe diretamente um único método $M_t^*$; somente $M_t^*$ é
decomposto.
Dada uma tarefa composta, o planner deverá:

\begin{enumerate}
\item identificar todos os métodos cujas pré-condições são satisfeitas;
\item expor uma representação estável desses métodos ao seletor;
\item aplicar uma máscara de validade durante treinamento e inferência;
\item solicitar ao MORL a escolha direta da alternativa preferida;
\item aplicar \emph{method fallback} somente se a decomposição escolhida falhar;
\item registrar métodos rejeitados, valores estimados e método escolhido.
\end{enumerate}

É importante distinguir inaplicabilidade de falha posterior.
A máscara HTN remove métodos que já violam restrições no estado corrente e esses
 métodos nunca chegam ao MORL. Se uma decomposição inicialmente aplicável não
 produz um plano executável após a expansão de suas subtarefas, o método
 selecionado é mascarado temporariamente e o MORL escolhe novamente entre as
 alternativas restantes.
 Essa recuperação não é uma busca DFS sobre uma lista ordenada: é excepcional e
normalmente o ciclo seleciona e decompõe apenas um método.
Quando o mundo muda durante a execução, o agente deve replanejar a partir do
estado observado, e não fazer backtracking sobre um estado simbólico obsoleto.
A separação entre planejamento inicial e replanejamento durante execução é
consistente com trabalhos recentes de replanning HTN sob mudanças de utilidade
e risco \cite{nabakowski2025}.

\begin{figure}[htbp]
\centering
\resizebox{0.92\textwidth}{!}{%
\begin{tikzpicture}[
    font=\small,
    >=Latex,
    node distance=10mm and 18mm,
    block/.style={
        rectangle, rounded corners=2mm, draw, thick, align=center,
        minimum width=35mm, minimum height=10mm
    },
    htn/.style={block, draw=htnblue, fill=htnblue!8},
    morl/.style={block, draw=morlred, fill=morlred!8},
    fallback/.style={block, draw=fallbackorange, fill=fallbackorange!10},
    runtime/.style={block, draw=envgray, fill=envgray!7},
    decision/.style={
        diamond, aspect=2.1, draw=black!70, fill=black!3, thick, align=center,
        inner sep=1.5pt
    },
    arrow/.style={->, thick}
]
    \node[htn] (compound) {Compound Task $C_t$};
    \node[htn, below=of compound] (filter) {HTN filters methods};
    \node[morl, below=of filter] (select) {MORL selects $M_t^*$};
    \node[htn, below=of select] (decompose) {Decompose only $M_t^*$};
    \node[decision, below=12mm of decompose] (success)
        {Decomposition\\successful?};
    \node[runtime, below left=15mm and 22mm of success] (execute)
        {Execute resulting\\primitive tasks};

    \node[fallback, below right=15mm and 22mm of success] (mask)
        {Mask failed $M_t^*$};
    \node[fallback, below=of mask] (remaining)
        {Updated feasible-method mask};
    \node[decision, below=of remaining] (left)
        {Any methods\\left?};
    \node[runtime, below right=13mm and 18mm of left] (failure)
        {Compound Task\\Failure};

    \node[decision, below=of execute] (invalid)
        {Runtime state\\invalidates plan?};
    \node[runtime, below left=14mm and 19mm of invalid] (continue)
        {Continue execution};
    \node[runtime, below right=14mm and 19mm of invalid] (replan)
        {Sense current $WS$\\and replan};

    \draw[arrow, htnblue] (compound) -- (filter);
    \draw[arrow, morlred] (filter) -- (select);
    \draw[arrow, htnblue] (select) -- (decompose);
    \draw[arrow] (decompose) -- (success);
    \draw[arrow] (success) -- node[above left] {yes} (execute);
    \draw[arrow, fallbackorange] (success) --
        node[above right] {no} (mask);
    \draw[arrow, fallbackorange] (mask) -- (remaining);
    \draw[arrow, fallbackorange] (remaining) -- (left);
    \draw[arrow, fallbackorange, rounded corners] (left.east)
        -- ++(26mm,0) |- node[pos=.22,right,align=left]
        {yes: fallback\\selection} (select.east);
    \draw[arrow] (left) -- node[above right] {no} (failure);

    \draw[arrow] (execute) -- (invalid);
    \draw[arrow] (invalid) -- node[above left] {no} (continue);
    \draw[arrow] (invalid) -- node[above right] {yes} (replan);
    \draw[arrow, envgray, dashed, rounded corners] (replan.east)
        -- ++(28mm,0) |- node[pos=.22,right,align=left]
        {current\\World State} (compound.east);
\end{tikzpicture}%
}
\caption{Separação entre o caminho online normal, a recuperação por
\emph{method fallback} e o replanejamento. O fallback só ocorre após falha de
decomposição e retorna ao MORL com uma máscara atualizada. Mudanças durante a
execução provocam replanejamento a partir do estado corrente, não backtracking
sobre um estado simbólico obsoleto.}
\label{fig:htn-morl-fallback}
\end{figure}


\section{Interface prevista}

A interface mínima entre os componentes deverá conter:

\begin{itemize}
\item estado ou representação de características $\phi(s_t)$;
\item identificador da tarefa composta $\tau_t$;
\item conjunto e máscara de métodos válidos;
\item vetor normalizado de preferências $\mathbf{w}_t$;
\item método selecionado $m_t$;
\item recompensa vetorial acumulada durante a decomposição;
\item duração da decisão e condição de término.
\end{itemize}

\section{Ambiente experimental alvo}

A próxima versão do ambiente deverá representar um cenário tático ou de
sobrevivência controlado.
Exemplos de tarefas compostas incluem defender uma posição, proteger um aliado,
obter acesso a uma área, neutralizar uma ameaça ou recuar sob recursos limitados.
Eventos estocásticos e mudanças de prioridade devem ser configuráveis e
reproduzíveis por sementes.

Uma evolução inspirada em ambientes multiobjetivo como
\textit{Resource Gathering} pode ser utilizada para validar recompensas
vetoriais e ferramentas de treinamento \cite{resourcegathering,mogymnasium,morlbaselines}.
Entretanto, o experimento principal deverá conter escolhas hierárquicas com
múltiplos métodos, pois a contribuição científica está na seleção online de
métodos HTN, e não apenas na navegação multiobjetivo.

Ambientes mais complexos, como Craftax, ficam reservados para validação futura
de escalabilidade e não fazem parte do escopo mínimo do mestrado
\cite{craftax2024,craftaxgithub}. Para padronizar a interação entre ambiente e
algoritmos, a implementação experimental deverá manter uma interface compatível
com Gymnasium \cite{gymnasium}.

\section{Instrumentação requerida}

Cada decisão deverá registrar estado, tarefa composta, métodos aplicáveis,
métodos rejeitados e respectivas pré-condições, vetor de preferências, vetor de
valor estimado por método, utilidade escalarizada, método escolhido, sequência
de subtarefas, recompensa vetorial, duração e resultado.
Esses registros sustentarão tanto a análise quantitativa quanto as explicações
qualitativas.
